\section{视频理解和大模型}

\subsection{为什么视频理解会走向 LLM}

视频模型下一步不只是分类器，而是可以被 prompt 的系统：输入视频和问题，输出描述、解释或结构化答案。课程最后提到的 Video-LLaVA、Video-ChatGPT、VideoLLaMA 3，都说明视频理解正在进入统一的多模态对话框架。

\begin{importantbox}{视频 LLM 的意义}
\begin{itemize}
    \item 它把视频理解从分类器推进到交互式系统。
\item 它把动作识别推进到语言解释和任务规划。
    \item 它也把视频理解和具身智能、助手系统更紧密地连起来。
\end{itemize}
\end{importantbox}

\subsection{这一讲真正想让你记住什么}

这节课表面上讲的是视频理解方法史，实际上讲的是一条更大的路径：从短 clip 到长视频，从外观到运动，从单模态到多模态，从分类到定位，从识别到问答，再从问答走向可操作的系统。

\begin{table}[H]
\centering
\begin{tabular}{p{3.1cm}p{3.4cm}p{4.8cm}}
\toprule
\textbf{阶段} & \textbf{核心问题} & \textbf{主要方法} \\
\midrule
短 clip 分类 & 动作是什么 & 单帧、late fusion、early fusion、3D CNN \\
运动分离 & 动作靠什么线索 & optical flow、two-stream \\
长时序建模 & 动作如何跨越更长时间 & CNN + RNN、recurrent CNN、attention \\
复杂任务 & 动作何时何地发生 & temporal localization、spatio-temporal detection \\
多模态与未来 & 还能利用什么信息 & 音视频融合、efficiency、LLM \\
\bottomrule
\end{tabular}
\caption{视频理解从局部动作识别一路扩展到多模态和 foundation model}
\end{table}

\begin{importantbox}{这节课最值得记住的几句话}
\begin{enumerate}
    \item 视频就是 \textbf{2D + time}，但时间维度会带来完全不同的建模挑战。
    \item 单帧 baseline 往往不差，但它忽略了真正的运动信息。
    \item 3D CNN 的本质是让卷积在时空体上滑动，而不是只在图像平面上滑动。
    \item Two-stream 说明了外观和运动可以分开建模，光流是经典运动表示。
    \item 长视频不能只靠短 clip；需要 RNN、attention 或更高效的时序聚合。
    \item 现实系统里，数据规模、采样策略和效率和模型结构同样重要。
\end{enumerate}
\end{importantbox}

